%% MoCafe v2.00 -- Physics and Algorithms
%% Author: Kwang-Il Seon (KASI)
\documentclass[english,a4paper]{kiseon_memo}
\synctex=1

\usepackage{booktabs}
\usepackage{amsmath}
\usepackage{amssymb}
\usepackage{graphicx}
\usepackage{microtype}

\newcommand{\kabs}{\kappa_{\mathrm{abs}}}
\newcommand{\kext}{\kappa_{\mathrm{ext}}}
\newcommand{\ksca}{\kappa_{\mathrm{sca}}}
\newcommand{\Cabs}{C_{\mathrm{abs}}}
\newcommand{\Cext}{C_{\mathrm{ext}}}
\newcommand{\Csca}{C_{\mathrm{sca}}}
\newcommand{\Jlam}{J_\lambda}
\newcommand{\khat}{\hat{\mathbf k}}
\newcommand{\nhat}{\hat{\mathbf n}}
\newcommand{\rvec}{\mathbf r}
\newcommand{\dd}{\mathrm{d}}
\newcommand{\code}[1]{\texttt{#1}}
\newcommand{\file}[1]{\texttt{#1}}
\newcommand{\pp}[1]{\texttt{par\%#1}}

\begin{document}

\title{MoCafe v2.00 --- Physics and Algorithms}
\author{Kwang-Il Seon (KASI)}
%--- compile-time date in ISO format (auto-updates on every build).
\date{\the\year-\ifnum\month<10 0\fi\the\month-\ifnum\day<10 0\fi\the\day}
\maketitle

\begin{abstract}
\noindent
This note documents the physics and the numerical algorithms behind MoCafe, the
parallel Monte-Carlo dust radiative-transfer code.  It develops the radiative
transfer problem and the Monte-Carlo solution (photon packets, free-path
sampling, forced first scattering, weighting, and Russian roulette), the
peeling-off (next-event) estimator, the scattering physics (Henyey--Greenstein
and the polarized Mueller-matrix path), panchromatic transport by grey
rescaling, the Lucy~(1999) path-length estimator for the radiation field, the
two dust-emission solvers (the Lucy iteration coupled to the \textsc{SEDust}
grain library with quantum-mechanical stochastic heating, and the
Bjorkman \& Wood immediate-reemission scheme), the Modified Random Walk, the
galaxy source models, the random-sampling routines (inversion, rejection, and
Walker's alias method), the continuous wavelength samplers of the source and
emission spectra, the quasi-random (Owen-scrambled Sobol) launch option,
and the HEALPix interior observer.  Derivations are
given at the level needed to read and modify the code; the companion user guide
covers installation, the input reference, and the outputs.
\end{abstract}

\tableofcontents
\clearpage

%=============================================================================
\section{The radiative transfer problem}
%=============================================================================

\subsection{The transfer equation}
The specific intensity $I_\lambda(\rvec,\nhat)$ of radiation of wavelength
$\lambda$ at position $\rvec$ travelling in direction $\nhat$ obeys the
radiative transfer equation
\begin{equation}
  \nhat\!\cdot\!\nabla I_\lambda
   = -\kext(\rvec,\lambda)\,I_\lambda
     + \kext(\rvec,\lambda)\,S_\lambda(\rvec,\nhat),
  \label{eq:rte}
\end{equation}
where $\kext=\ksca+\kabs$ is the extinction coefficient (per unit length),
$\ksca$ and $\kabs$ the scattering and absorption coefficients, and $S_\lambda$
the source function.  For a dusty medium the source function has a scattering
part and a thermal-emission part,
\begin{equation}
  S_\lambda(\rvec,\nhat)
   = \frac{\ksca}{\kext}\!\int\! \frac{\dd\Omega'}{4\pi}\,
       \Phi(\nhat,\nhat')\,I_\lambda(\rvec,\nhat')
   + \frac{\kabs}{\kext}\, B_\lambda\!\big(T_{\rm d}(\rvec)\big),
  \label{eq:source}
\end{equation}
with $\Phi$ the (normalized) scattering phase function, $a\equiv\ksca/\kext$ the
single-scattering albedo, and $B_\lambda(T_{\rm d})$ the thermal emission of dust
at temperature $T_{\rm d}$.  Defining the optical depth along a ray by
$\dd\tau_\lambda=\kext\,\dd\ell$, the formal solution of Eq.~\eqref{eq:rte} is
\begin{equation}
  I_\lambda(\tau_\lambda)= I_\lambda(0)\,e^{-\tau_\lambda}
     + \int_0^{\tau_\lambda}\!\! S_\lambda(\tau')\,e^{-(\tau_\lambda-\tau')}\,\dd\tau',
  \label{eq:formal}
\end{equation}
i.e.\ attenuated incident light plus the attenuated contribution of every source
point along the ray.

\subsection{Why Monte Carlo}
Equations~\eqref{eq:rte}--\eqref{eq:source} couple every direction and position
(scattering redistributes light in angle; thermal emission couples wavelengths
through the temperature).  Rather than discretizing the integro-differential
operator, MoCafe evaluates the physical content of Eq.~\eqref{eq:formal}
\emph{statistically}: it follows a large number of photon packets whose random
histories --- emission, propagation over an exponentially distributed optical
depth, scattering into a new direction, absorption --- reproduce the transfer
equation in expectation.  Integrals over the phase function, the size
distribution, and the wavelength grid become sums over sampled events, so the
cost scales with the number of packets rather than with the dimensionality of
the problem, and complex three-dimensional geometries (Cartesian, clumpy, AMR)
are handled by the same estimator.

%=============================================================================
\section{Monte-Carlo photon transport}
%=============================================================================

\subsection{Photon packets and weights}
MoCafe transports photon \emph{packets}: weighted bundles of monochromatic
photons.  A packet carries a position $\rvec$, a propagation direction $\khat$,
a dimensionless weight $w$, a wavelength $\lambda$, a Stokes vector
$(I,Q,U,V)$ when polarization is on, and---in SED/emission mode---a physical
luminosity $L_{\rm p}$ [erg\,s$^{-1}$] so that energy budgets are absolute.  The
weight absorbs all the variance-reduction bookkeeping (forced scattering,
albedo, scan reweighting); the packet contributes $w\,L_{\rm p}$ wherever it is
tallied.

\subsection{Sampling the free path}
The probability that a packet reaches optical depth $\tau$ without interacting is
$e^{-\tau}$, so the interaction optical depth has probability density
$p(\tau)=e^{-\tau}$ and cumulative $P(\tau)=1-e^{-\tau}$.  Inverting
$P(\tau)=\xi$ for a uniform deviate $\xi\in[0,1)$ gives the standard free-path
sample
\begin{equation}
  \tau = -\ln(1-\xi).
  \label{eq:freepath}
\end{equation}
The sampled $\tau$ is converted to a physical displacement by integrating
$\kext$ along the ray until the accumulated optical depth equals $\tau$.  The
three media do this with a shared interface: Amanatides--Woo voxel traversal on
the Cartesian grid, a ray--sphere intersection followed by a DDA walk over the
clump acceleration grid in the clumpy medium, and an octree neighbour walk on the
AMR grid.  Because scattering and peeling-off call these tracers through a
procedure pointer, they are written once and are grid-agnostic.

\subsection{Forced first scattering and reduced weight}
In an optically thin medium most packets would escape without interacting,
which wastes samples.  Setting \pp{use\_reduced\_wgt} forces the first
interaction to occur inside the grid.  Let $\tau_0$ be the total optical depth
from the birth point to the far edge along $\khat$.  Restricting the free-path
distribution to $[0,\tau_0]$ and renormalizing gives the truncated density
$p(\tau)=e^{-\tau}/(1-e^{-\tau_0})$, whose inverse is
\begin{equation}
  \tau = -\ln\!\big[\,1-\xi\,(1-e^{-\tau_0})\big],\qquad \tau\in[0,\tau_0].
\end{equation}
To keep the estimator unbiased the packet weight is multiplied by the
probability that it would have interacted at all,
\begin{equation}
  w \leftarrow w\,(1-e^{-\tau_0}),
\end{equation}
and the complementary fraction $e^{-\tau_0}$ --- the probability of escaping
without scattering --- is carried by the direct (unscattered) peel-off toward
each observer, so no energy is lost.  This is why the $n_{\rm scatt}=0$ step is
special-cased in the transport loop.

\subsection{Absorption, weighting, and Russian roulette}
At an interaction MoCafe uses \emph{albedo weighting} rather than a stochastic
absorb/scatter decision: the packet always scatters, and the absorbed fraction
is accounted continuously (the scattered peel carries a factor $a$; the absorbed
energy is deposited into the radiation-field tally, \S\ref{sec:jlam}).  This
removes the variance of the absorb/scatter coin.  A packet whose weight has
decayed below a cutoff is subjected to Russian roulette: with survival
probability $q$ it continues with weight $w/q$, otherwise it is terminated.
Roulette keeps the estimator unbiased while preventing unbounded tracking of
negligible packets, and is used both in the transport loop and inside the
Modified Random Walk (\S\ref{sec:mrw}).

%=============================================================================
\section{The peeling-off (next-event) estimator}
%=============================================================================

\subsection{Derivation}
Waiting for packets to reach a distant observer by chance is hopelessly
inefficient.  Instead, at every emission and scattering site MoCafe adds the
\emph{expected} contribution the packet would make toward each observer --- the
peeling-off or next-event estimator.  Consider a scattering at position $\rvec$
with incoming direction $\khat$.  The joint probability that the packet scatters
into the small solid angle $\dd\Omega_o$ about the observer direction $\nhat_o$
\emph{and} then escapes to the observer without further interaction is
\begin{equation}
  a\,\frac{\Phi(\khat\!\cdot\!\nhat_o)}{4\pi}\,\dd\Omega_o \times e^{-\tau_o},
\end{equation}
where $a$ is the albedo, $\Phi$ the normalized phase function, and $\tau_o$ the
optical depth from $\rvec$ to the observer along $\nhat_o$.  For an observer of
collecting area $\dd A$ at distance $d_o$, $\dd\Omega_o=\dd A/d_o^2$, so the
image-plane contribution per unit area is
\begin{equation}
  \Delta I_o = \frac{w\,a\,\Phi(\khat\!\cdot\!\nhat_o)}{4\pi}\,
               \frac{e^{-\tau_o}}{d_o^{\,2}}.
  \label{eq:peel}
\end{equation}
The $1/d_o^2$ is the geometric dilution, $e^{-\tau_o}$ the attenuation to the
observer (computed by ray-tracing to the grid edge), and the phase-function
factor the probability of redirection.  Summing Eq.~\eqref{eq:peel} over all
scattering events of all packets builds the scattered image with far lower
variance than collecting escaping packets, because \emph{every} event
contributes to \emph{every} observer.

\subsection{Direct and scattered peels}
The unscattered (direct) light is peeled separately: a source packet contributes
its birth luminosity attenuated by the optical depth from the source to the
observer, forming the \code{Direct} image (and \code{Direct0}, the unattenuated
version, when requested).  Subsequent scatterings feed the \code{Scattered}
image through Eq.~\eqref{eq:peel}.  For an internal all-sky observer
(\S\ref{sec:allsky}) the same expected contribution is binned into the HEALPix
pixel of the arrival direction instead of a tangent-plane pixel.  All peels go
through the \code{raytrace\_to\_edge} procedure pointer, so the Cartesian,
clumpy, and AMR media share identical scattering and peeling code; the default
binding keeps the Cartesian output bit-identical across refactors.

%=============================================================================
\section{Scattering}
%=============================================================================

\subsection{Henyey--Greenstein}
The default phase function is Henyey--Greenstein with asymmetry parameter
$g=\langle\cos\theta\rangle$,
\begin{equation}
  \Phi_{\rm HG}(\mu) = \frac{1-g^2}{(1+g^2-2g\mu)^{3/2}},\qquad \mu=\cos\theta,
\end{equation}
normalized so that $\tfrac12\!\int_{-1}^{1}\Phi_{\rm HG}\,\dd\mu = 1$.  Its
cumulative distribution inverts in closed form, so the scattering cosine is
sampled directly from a uniform deviate $\xi$:
\begin{equation}
  \mu = \frac{1}{2g}\left[1+g^2-\left(\frac{1-g^2}{1-g+2g\xi}\right)^{2}\right]
        \quad (g\neq0),
\end{equation}
with the isotropic limit $\mu=2\xi-1$ for $g=0$.  The azimuth is uniform, and
the new direction is built by rotating $\khat$ by $(\theta,\varphi)$.

\subsection{Polarization: Stokes vectors and the Mueller matrix}
When a scattering-matrix table is supplied (\pp{scatt\_mat\_file}) the code
carries the Stokes vector $\mathbf S=(I,Q,U,V)^{\!\top}$ and replaces the scalar
phase function by the $4\times4$ Mueller matrix $\mathbf M(\theta)$ of the
grains.  A scattering event proceeds in three steps: (i) rotate $\mathbf S$ into
the scattering plane by the Stokes rotation $\mathbf R(\psi_1)$; (ii) apply the
Mueller matrix, $\mathbf S'=\mathbf M(\theta)\,\mathbf R(\psi_1)\,\mathbf S$;
(iii) rotate the result into the new meridian frame by $\mathbf R(\psi_2)$.  The
scattering angle is drawn from the tabulated $S_{11}(\theta)$ (the intensity
phase function) with an alias table (\S\ref{sec:alias}) for $\mathcal O(1)$
sampling; because $S_{11}$ already encodes the angular dependence, the azimuth is
drawn from the $S_{12}$-modulated distribution to respect the incident
polarization.  The peel-off of Eq.~\eqref{eq:peel} then uses the
matrix-weighted phase function evaluated for the observer direction, so the
polarized images ($Q,U,V$ planes) are next-event estimates just like $I$.  With
Henyey--Greenstein, \pp{use\_stokes} is forced off.

%=============================================================================
\section{Panchromatic (SED) transport}\label{sec:sed}
%=============================================================================

\subsection{Grey rescaling}
In SED mode a packet is assigned a wavelength at birth and keeps it until
absorption.  Storing a separate opacity grid per wavelength would multiply the
memory by the number of bins; instead MoCafe stores the density structure --- and
hence the optical depth --- at a single reference wavelength $\lambda_{\rm ref}$
and rescales each packet by the wavelength-dependent opacity ratio
(Jonsson~2006),
\begin{equation}
  s_{\rm ext}(\lambda) = \frac{\kext(\lambda)}{\kext(\lambda_{\rm ref})},
  \qquad
  \tau(\lambda) = s_{\rm ext}(\lambda)\,\tau_{\rm ref}.
  \label{eq:grey}
\end{equation}
The ray tracers integrate $\tau_{\rm ref}$ and multiply by the packet's
$s_{\rm ext}$; the free-path draw of Eq.~\eqref{eq:freepath} is done in
$\tau_{\rm ref}=\tau(\lambda)/s_{\rm ext}$.  Because $s_{\rm ext}$ is a single
scalar per packet, all three tracers are shared with the monochromatic path and
the memory does not grow with the number of wavelengths.

\subsection{Where the cross sections come from}\label{sec:optics}
$\Cext(\lambda)$, the albedo $a(\lambda)$ and the asymmetry $g(\lambda)$ are the
size-integrated optics of the grain model named by \pp{dust\_model},
\begin{equation}
  \Cext(\lambda) = \sum_{p}\sum_{a} \dd n_p(a)
     \big[\,C_{\rm abs}^{\,p}(\lambda,a) + C_{\rm sca}^{\,p}(\lambda,a)\,\big],
  \qquad
  a(\lambda) = \frac{C_{\rm sca}}{\Cext},
  \label{eq:sizeint}
\end{equation}
\begin{equation}
  g(\lambda) = \frac{\sum_p\sum_a \dd n_p(a)\,C_{\rm sca}^{\,p}(\lambda,a)\,
                     g^{\,p}(\lambda,a)}
                    {\sum_p\sum_a \dd n_p(a)\,C_{\rm sca}^{\,p}(\lambda,a)},
  \label{eq:gbar}
\end{equation}
summed over every grain population $p$ of the model with $\dd n_p(a)$ its number
per H atom in the size bin at $a$.  \textsc{SEDust} evaluates both from the same
model object that later supplies the emission, so the energy a cell absorbs is
set by the $C_{\rm abs}$ of the grains that then radiate it away; nothing in the
energy balance depends on two independently specified dust models agreeing.

Each model supplies its scattering optics in the way that model is defined:
\texttt{astrodust} from its random-orientation T-matrix $Q$ table; \texttt{dl07}
and \texttt{mrn} from Mie theory on the D03 astrosilicate and graphite
dielectric functions, the DL07 PAH component entering through absorption alone
because a PAH scatters negligibly and MRN having no PAHs at all;
\texttt{zubko} and a file-defined model from the $Q_{\rm sca}$ and
$\langle\cos\rangle$ columns of their own tables.

An explicit table \pp{kext\_file} ($\lambda$, $a$, $\langle\cos\rangle$,
$\Cext/$H) overrides this.  It is then the user's responsibility that the table
came from the same model as \pp{dust\_model}: MoCafe cannot check it, and a
mismatch breaks the balance above without any symptom in the transport.

\subsection{Wavelength sampling and packet luminosity}\label{sec:wavesamp}
The source spectrum $L_\lambda$ --- a Planck function at \pp{tstar}, a tabulated
\pp{source\_spectrum}, or a spectrum for each component in the multi-source case
--- is carried as a piecewise-linear density on a node set that holds both the
tabulated wavelengths of the spectrum and every edge of the transfer grid; a
blackbody is refined adaptively to a relative tolerance before entering the same
representation.  The cumulative distribution of that density integrates in closed
form, so a packet's birth wavelength follows by inverting it,
\begin{equation}
  \int_{\lambda_{\min}}^{\lambda_{\rm p}} L_\lambda\,{\rm d}\lambda
  \;=\; \xi \int_{\lambda_{\min}}^{\lambda_{\max}} L_\lambda\,{\rm d}\lambda ,
  \qquad \xi \sim U(0,1),
  \label{eq:wavedraw}
\end{equation}
which inside the bracketing interval is a quadratic solved in closed form.  The
wavelength is therefore \emph{continuous}: drawing a bin index and placing the
packet at the bin center would freeze $s_{\rm ext}$, $a$ and $g$ at their
bin-center values and erase whatever spectral structure the bin contains.  The
packet still carries the index of the bin its wavelength falls in, because the
output images and the radiation-field tally are binned in wavelength, but no
optical quantity is read at a bin center.

Because the transfer bin edges are nodes of the density, the luminosity of each
bin is an exact difference of the cumulative distribution rather than the
midpoint estimate $L_\lambda(\lambda_i)\,\Delta\lambda_i$; on a coarse grid the
two differ substantially (for a blackbody on four bins per decade the midpoint
rule errs by tens of percent).  Each packet carries a physical luminosity
$L_{\rm p}=L_{\rm tot}/N_{\rm phot}$, so that the sum of packet luminosities
reproduces the source luminosity and the downstream emission budget is absolute
(erg\,s$^{-1}$).  Multiple stellar populations (\S\ref{sec:galaxy}) are mixed by
drawing the emitting component in proportion to its luminosity.

The source spectrum may be supplied in two ways.  As a \emph{shape}, the file is
renormalized and the absolute scale comes from \pp{luminosity}.  As an
\emph{absolute} (physical) spectrum, the source luminosity is instead the integral
of the tabulated $L_\lambda$ over the wavelength grid,
\begin{equation}
  L_{\rm tot}=\int L_\lambda\,\dd\lambda ,
\end{equation}
and inputs tabulated in frequency or energy are converted with the Jacobians
\begin{equation}
  L_\lambda = L_\nu\,\frac{c}{\lambda^2},
  \qquad
  L_\lambda = L_E\,\frac{hc}{\lambda^2}.
\end{equation}
An isotropic external radiation field is specified instead by its mean intensity
$J$; the luminosity it injects across an illuminated boundary of area
$A_{\rm surface}$ is
\begin{equation}
  L_{\rm tot}=\pi\,J\,A_{\rm surface},
  \label{eq:extlum}
\end{equation}
after which the packets are drawn and weighted exactly as for an internal source.

\paragraph{Composed sources.}  An internal source (or sources) and an external
isotropic field may coexist in one run.  Each packet is then assigned an origin
--- internal or external --- with probability $L_{\rm int}:L_{\rm ext}$, and all
packets share the luminosity
\begin{equation}
  L_{\rm p}=\frac{L_{\rm tot}}{N_{\rm phot}},\qquad
  L_{\rm tot}=L_{\rm int}+L_{\rm ext}
            =\sum_i L_i + \pi\,J\,A_{\rm surface},
  \label{eq:composelum}
\end{equation}
so the absolute normalization is the sum of the internal luminosity
$L_{\rm int}=\sum_i L_i$ (Eq.~\ref{eq:extlum} with $A_{\rm surface}$ the
illuminated boundary) and the injected external luminosity.  Because energy is
additive, the composed image equals the internal-only image plus the
external-only image, and the output is normalized to $L_{\rm tot}$.

\subsection{Quasi-random photon launching}\label{sec:qmc}
The launch of a packet is a fixed, low-dimensional integral over the unit
hypercube: the origin (internal or external), the emitting component, the
wavelength, the emission (or incidence) direction, and, for external
illumination, the entry point on the boundary.  With
\pp{launch\_sequence}${}=$\code{'sobol'} these coordinates are taken from an
Owen-scrambled Sobol sequence instead of independent pseudo-random numbers.  A
low-discrepancy point set fills the launch cube more uniformly: for a smooth
$d$-dimensional integrand the quadrature error can approach
$N^{-1}(\log N)^{d}$ instead of the Monte-Carlo $N^{-1/2}$, and even for the
discontinuous transport integrand the stratification is retained where it
matters most.  In particular, inverting one stratified coordinate through the
monotonic luminosity CDF makes the packet count in a wavelength bin of
probability $p$ deviate from $pN$ by order one rather than binomially.  Both
launch paths invert that CDF (\S\ref{sec:wavesamp}); nothing may permute the
unit interval on this coordinate, since the permutation would destroy the
monotonicity the stratification rests on.

The design is deliberately hybrid.  Only the launch coordinates --- a fixed
seven-dimensional layout whose meanings never depend on the source
configuration --- come from the scrambled net, indexed by the global photon
number carried as \code{photon\%id}; every conditional, variable-length draw
after launch (forced first scattering, the phase function, free paths) stays
on the Mersenne Twister, so the transport estimator is unchanged and
unbiased.  Owen scrambling (nested-uniform hashing) preserves the net's
equidistribution while making replicates with different \pp{qmc\_seed} values
statistically independent, which restores an ordinary error estimate from the
replicate-to-replicate scatter.  Because the launch set is indexed globally,
it is independent of the MPI task count; the direct images of internal
sources, which depend on the launch variables alone, are then bit-identical
for any number of ranks.  The scattered field keeps its $N^{-1/2}$ behavior:
the benefit concentrates in the wavelength allocation, the direct images, and
the entry-point uniformity of external illumination.

%=============================================================================
\section{The radiation field: Lucy path-length estimator}\label{sec:jlam}
%=============================================================================

\subsection{Mean intensity}
Dust heating is set by the mean intensity (zeroth angular moment of $I_\lambda$)
in each cell,
\begin{equation}
  \Jlam(\rvec)=\frac{1}{4\pi}\!\int I_\lambda(\rvec,\nhat)\,\dd\Omega .
\end{equation}
MoCafe estimates $\Jlam$ with the Lucy~(1999) \emph{path-length} estimator.  The
energy density is $u_\lambda = (4\pi/c)\Jlam$, and the total path length that
radiation traverses in a volume $V$ per unit time equals $c$ times the number of
photons in $V$; accumulating a contribution from every segment a packet takes
through the cell therefore estimates $\Jlam$ directly:
\begin{equation}
  \Jlam = \frac{1}{4\pi\,V\,\Delta\lambda\,d_{\rm cm}^2}
          \sum_{\text{segments in cell}} L_{\rm p}\, w\, \ell,
  \label{eq:jlucy}
\end{equation}
where $\ell$ is the segment length, $V$ the cell volume, $\Delta\lambda$ the bin
width, and $d_{\rm cm}$ the length unit in cm.

\subsection{Path-length versus collision estimator}
A collision (or absorption) estimator would tally only at discrete interaction
points, so a cell with few interactions receives a noisy, or zero, estimate.
The path-length estimator instead credits every fly-through, whether or not the
packet interacts, so even optically thin cells --- where interactions are rare
but many packets pass through --- get a smooth, low-variance $\Jlam$.  This is
decisive for dust heating, where the illuminating field must be known in cells
that are optically thin to the starlight.

\subsection{Absorbed power and the analytic first flight}
The power absorbed in a cell follows from $\Jlam$ and the absorption opacity,
\begin{equation}
  P_{\rm abs} = 4\pi V \!\int \kabs(\lambda)\,\Jlam \, \dd\lambda .
  \label{eq:pabs}
\end{equation}
This integral is accumulated as the packets fly, not assembled afterwards from
the binned $\Jlam$: each segment contributes its path-length times
$\Cabs(\lambda_{\rm p})$ at the packet's \emph{own} wavelength.  Forming it
instead from $\Jlam$ would have to use one cross section for a whole bin, and
because $\Cabs$ varies across a bin that biases the absorbed --- hence the
emitted --- power of every cell by a few percent, which propagates straight into
the dust temperature.  The pathlength and absorption tallies are accumulated
together, so this costs one array over cells and no extra traversal.

The forced first flight (\S\ref{sec:sed}) has a biased free-path distribution, so
that segment is \emph{not} tallied by sampling; instead its exact expectation is
added analytically along the ray to the edge,
\begin{equation}
  \Jlam \mathrel{+}= \frac{w\,L_{\rm p}}{4\pi V\Delta\lambda\,d_{\rm cm}^2}
     \int e^{-\tau_\lambda(\ell)}\,\dd\ell
   = \frac{w\,L_{\rm p}}{\dots}\,
     \frac{e^{-s\,\tau_{\rm in}}-e^{-s\,\tau_{\rm out}}}{\kext s},
\end{equation}
a zero-variance direct component.  Flights with $n_{\rm scatt}\ge1$ sample the
ordinary exponential free path and are tallied along the walked segments.  The
$\code{\_jlam}$ output reports the pathlength-tallied absorbed power
(\code{EABS\_A}) against an independent event counter (\code{EABS\_B}) as an
energy-conservation check.

%=============================================================================
\section{Dust thermal emission I: Lucy + SEDust}\label{sec:lucy}
%=============================================================================

The radiation field $\Jlam$ from Eq.~\eqref{eq:jlucy} is handed cell by cell to
the \textsc{SEDust} grain library, which returns the emission spectrum for the
local field.  Below, $\Cabs(a,\lambda)=\pi a^2 Q_{\rm abs}(a,\lambda)$ is the
absorption cross section of a grain of radius $a$.

\subsection{The radiation field handed to SEDust}\label{sec:jhandoff}
The two codes work on different wavelength grids: the transfer grid is the
\pp{nlambda} log-spaced bins, while the \textsc{SEDust} grid $\{\lambda_j\}$ is
fixed by the optical-constant tables of the grain model (1129 points for
astrodust).  Neither is a refinement of the other, and their relative resolution
is a user choice, so $\Jlam$ has to be carried across without assuming which is
finer.

What the tally holds is not a function but a set of \emph{bin energies}: bin $i$,
spanning $[\lambda_i,\lambda_{i+1}]$, has received
\begin{equation}
  E_i \;=\; \int_{\lambda_i}^{\lambda_{i+1}}\! \Jlam \,\dd\lambda
  \;=\; \bar J_i\,\Delta\lambda_i ,
  \label{eq:binenergy}
\end{equation}
with $\bar J_i$ the bin average.  A plain resample reads $\bar J_i$ as though it
were the point value at the bin center and interpolates through those points.
The interpolant $\tilde J_\lambda$ that results is smooth, but nothing ties it to
Eq.~\eqref{eq:binenergy}: its own bin integrals
$\int_{\rm bin}\tilde J_\lambda\,\dd\lambda$ differ from $E_i$ by
$13\%$ in the mean over a grid running from the ultraviolet to the
submillimetre, and the total by $1.7\%$.  Handing that to \textsc{SEDust} would
have the grains heated by a spectrum the transport never produced.

The shape of the interpolant is worth keeping --- it is smooth, positive, and
carries the slope information the neighbouring bins imply --- so it is kept and
only its normalization is corrected, bin by bin:
\begin{equation}
  \Jlam \;=\; \tilde J_\lambda \;
     \frac{E_i}{\displaystyle\int_{\lambda_i}^{\lambda_{i+1}}\!\tilde J_\lambda\,\dd\lambda},
  \qquad \lambda\in[\lambda_i,\lambda_{i+1}] .
  \label{eq:jrescale}
\end{equation}
Every bin then carries exactly the energy the packets deposited there.  Measured
on the same grid, the per-bin discrepancy falls from $13\%$ to $1.1\%$ in the
mean and the total from $1.7\%$ to $0.08\%$.  The residual is structural rather
than a loose end: Eq.~\eqref{eq:jrescale} applies a different factor on either
side of a bin edge, so $\Jlam$ is discontinuous there, and any quadrature that
interpolates across the edge --- \textsc{SEDust}'s own included --- mixes the two
factors.  Chasing it further is not worthwhile, since what remains sits below the
shot noise of the per-bin tally itself.

A staircase, $\Jlam=\bar J_i$ across each bin, would be exactly conservative
without any correction; it is rejected because it replaces one discontinuity per
bin edge with a spectrum that has no slope at all, and \textsc{SEDust} integrates
$\Cabs(\lambda)\Jlam$ over a grid fine enough to feel the difference.

A bin no packet ever reached is left empty rather than filled in from its
neighbours, since interpolating would place radiation at a wavelength the
transport says carries none.  The two reasons a bin can be empty --- radiation
genuinely absent, or present but never sampled --- are not distinguishable from
the tally, so this is a deliberate choice and it has a cost: a bin left empty by
shot noise loses its heating and biases the cell cool.  That bites where the
statistics are thin (deep in optically thick regions, far from the sources, or at
a low photon count) and worst in the ultraviolet, which both empties first and
drives the stochastic heating.  The remedy there is more packets, not
interpolation.

\subsection{Grain energy balance}
A grain absorbs starlight at the rate
$\dot E_{\rm abs}=\int 4\pi\,\Cabs(a,\lambda)\,\Jlam\,\dd\lambda$ and radiates at
$\dot E_{\rm emit}(T)=\int 4\pi\,\Cabs(a,\lambda)\,B_\lambda(T)\,\dd\lambda$.  A
\emph{large} grain reaches a steady temperature $T_{\rm eq}$ where
$\dot E_{\rm emit}(T_{\rm eq})=\dot E_{\rm abs}$; because the absorbed energy per
photon is tiny compared with the grain heat capacity, its temperature barely
fluctuates.  A \emph{small} grain (or a PAH molecule), by contrast, has so small
a heat capacity that a single ultraviolet photon can raise its temperature by
hundreds of kelvin, after which it cools before the next hit.  Its temperature
is therefore a random variable, and the emission must be integrated over its
probability distribution.

\subsection{Enthalpy and heat capacity}
The internal energy (enthalpy) $U(T)$ of a grain is obtained from the Debye
model of the lattice vibrations (Draine \& Li 2001): silicate and graphite use
two- and three-dimensional Debye integrals with the appropriate Debye
temperatures, and PAHs add discrete C--H and C--C mode contributions
(\file{enthalpy\_v2.f90}, \file{enthalpy\_astrodust.f90}).  Inverting $U(T)$
gives the temperature reached after absorbing a given energy, which is the input
to the stochastic-heating solver.

\subsection{Stochastic heating: the temperature probability distribution}
\textsc{SEDust} solves for the steady-state probability distribution $P_j$ over
a set of discrete enthalpy bins $\{U_j\}$ ($j=1\ldots N$, default $N=200$) using
Draine's thermal-discrete method (\file{stoch\_qm.f90}; Guhathakurta \&
Draine~1989; Draine \& Li~2001).  Transitions between bins are collected into a
rate matrix $A_{fi}$ (the rate of going from bin $i$ to bin $f$):
\begin{itemize}
\item \textbf{Heating} ($f>i$): absorption of a photon of energy
      $U_f-U_i=hc/\lambda$ carries the grain upward, at rate
      \begin{equation}
        A_{fi} = \frac{4\pi\,\Jlam}{hc/\lambda}\,\Cabs(a,\lambda)\,
                 \frac{\dd\lambda}{\dd U_f},
      \end{equation}
      i.e.\ proportional to the photon arrival rate at the wavelength that
      bridges the two bins.
\item \textbf{Cooling} ($f<i$): thermal emission removes energy.  In the
      thermal-discrete treatment (\code{dbdis}, the production setting) all
      downward transitions are retained, with rates set by the Planck emission
      at the temperature of the upper bin; the continuous-cooling approximation
      (\code{dbcon}) collapses these into nearest-neighbour ($f=i-1$)
      transitions, recovering the Guhathakurta--Draine temperature-space
      recursion.
\end{itemize}
Conservation of probability fixes the diagonal, and the steady state satisfies
\begin{equation}
  \sum_{i\neq j} A_{ji}\,P_i = P_j \sum_{i\neq j} A_{ij},
  \qquad \sum_j P_j = 1,
\end{equation}
solved as a sparse linear system with a bi-conjugate-gradient solver.  The
result $P_j$ is the probability that the grain sits in enthalpy bin $j$, i.e.\
at temperature $T_j=T(U_j)$.

\subsection{Equilibrium limit}
For large grains the distribution $P_j$ collapses to a narrow spike at
$T_{\rm eq}$, and the solve reduces to the energy-balance root
$\dot E_{\rm emit}(T_{\rm eq})=\dot E_{\rm abs}$.  \textsc{SEDust} detects this
regime and can use the equilibrium temperature directly, which is exact for the
big grains that dominate the far-infrared and much cheaper than the full solve.

\subsection{Size distribution and PAHs}
The cell emissivity integrates the single-grain emission over the grain-size
distribution $\dd n/\dd a$ of each material (silicate, carbonaceous, PAH),
\begin{equation}
  \varepsilon_\lambda = \sum_{\rm species}\int \dd a\;\frac{\dd n}{\dd a}\;
     4\pi\,\Cabs(a,\lambda)\int \dd T\; P(a,T)\, B_\lambda(T),
  \label{eq:emis}
\end{equation}
where $P(a,T)$ is the stochastic distribution above (a delta function at
$T_{\rm eq}$ for large grains).  The PAH population produces the mid-infrared
band spectrum (3.3, 6.2, 7.7, 8.6, 11.3\,$\mu$m); the smallest grains produce
the smoothly rising continuum from single-photon heating.  The size
distributions and optical constants come from the selected grain model.

\subsection{Emission spectrum and energy normalization}
The emission in each cell is renormalized so that its wavelength integral equals
the locally absorbed power of Eq.~\eqref{eq:pabs},
\begin{equation}
  \int \varepsilon_\lambda^{\rm norm}\,\dd\lambda = P_{\rm abs},
\end{equation}
which makes energy conservation \emph{exact} regardless of the internal
\textsc{SEDust} normalization.

The energy of each transfer bin is the integral of the emission spectrum over
that bin, taken on the \textsc{SEDust} wavelength grid, which is far finer than
the transfer grid and resolves the aromatic features and the Wien flank that lie
inside a bin.  Reading the spectrum once at the bin center and multiplying by
the bin width instead --- a midpoint rule --- misses that structure by several
percent in the mean and by tens of percent in the worst bin.  A two-point rule
on the bin edges is no remedy: the midpoint rule already carries half the error
of any endpoint pair ($-h^3f''/24$ against $+h^3f''/12$), so the only way past
it is to integrate the resolved spectrum, which is what is done.

Those bin energies give the intrinsic dust SED directly.  For the emergent SED
and the pixel-resolved image a second batch of packets is emitted, drawing the
cell from the emissivity and the bin from the same energies; the wavelength
inside the bin then follows from reading $\lambda\varepsilon_\lambda$ as a power
law between the two bin edges and inverting it, which for
$q=\lambda_{i+1}\varepsilon_{i+1}/\lambda_i\varepsilon_i$ and
$x=\ln(\lambda/\lambda_i)/\ln(\lambda_{i+1}/\lambda_i)$ is
\begin{equation}
  x = \frac{\ln\!\left[1+\xi\,(q-1)\right]}{\ln q},
  \qquad \xi \sim U(0,1).
  \label{eq:powerlawbin}
\end{equation}
Only the ratio $q$ enters, so the draw never disturbs the bin energies.  The
leftover of the uniform that selected the bin is itself uniform inside it and is
reused as $\xi$, so no extra random number is drawn and a stratified
quasi-random coordinate stays stratified.

\subsection{Speed options}
The exact stochastic solve costs $\mathcal O(0.1\,\mathrm{s})$ per cell.  Two
approximations trade accuracy for speed:
\begin{description}
\item[\pp{dust\_single\_teq}] replaces $P(a,T)$ by a single mixture equilibrium
  temperature.  It keeps the total power but drops the PAH and stochastic
  features (no mid-infrared bands), and is intended for optically thin,
  far-infrared-dominated problems.
\item[\pp{dust\_fast\_table}] exploits that, to good approximation, the emission
  spectrum depends on the field only through its intensity
  $U=\int\Jlam\,\dd\lambda/\int J_\lambda^{\rm ref}\,\dd\lambda$ and shape.  A
  reference emission table is built over a grid of $U$ once, and each cell is
  interpolated in $\log U$.  This preserves the PAH/stochastic features at
  $\sim1\%$ SED-shape error, with the largest deviations on the far-infrared
  slope.
\end{description}

\subsection{Self-absorption iteration}
At high infrared optical depth the dust re-absorbs its own emission.  With
\pp{dust\_niter}$>1$ MoCafe transports the emitted packets, tallies their
contribution to the field, adds it to the stellar field
$\Jlam\to\Jlam^{\star}+\Jlam^{\rm dust}$, and recomputes the emission, iterating
until the total emitted luminosity converges to \pp{dust\_tol}.  Writing $f$ for
the fraction of dust emission that is re-absorbed, the emitted luminosity obeys
\begin{equation}
  L^{(k)} = L_{\rm abs}^{\star} + f\,L^{(k-1)}
  \;\longrightarrow\;
  \frac{L_{\rm abs}^{\star}}{1-f},
\end{equation}
a geometric series.  At $\tau_V=100$ this converges in about ten iterations, and
only the converged run conserves energy: a single pass loses the re-absorbed
dust emission (only $\sim37\%$ of the starlight escapes), whereas the converged
run returns $99.8\%$.  Iteration is therefore \emph{required} for energy
conservation at high infrared optical depth.

%=============================================================================
\section{Dust thermal emission II: Bjorkman \& Wood}\label{sec:bw}
%=============================================================================

The Bjorkman \& Wood~(2001) scheme uses fixed-energy packets and a mixture mean
opacity, and needs no radiation-field storage and no iteration.  When a packet is
absorbed in a cell, the cell's cumulative absorbed energy increases, raising its
equilibrium temperature from $T$ to $T+\Delta T$; the packet is re-emitted
\emph{immediately} with a wavelength drawn from the temperature-correction
spectrum.  The construction is as follows.  A cell that has already emitted a
spectrum $\propto\kabs B_\lambda(T)$ must, after the temperature increment,
radiate $\propto\kabs B_\lambda(T+\Delta T)$.  The \emph{additional} photons
needed to correct the spectrum are therefore distributed as the difference,
\begin{equation}
  \frac{\dd P}{\dd\lambda}\;\propto\;
     \kabs(\lambda)\,\big[B_\lambda(T{+}\Delta T)-B_\lambda(T)\big]
  \;\approx\;
     \kabs(\lambda)\,\frac{\partial B_\lambda(T)}{\partial T},
  \label{eq:bwcorr}
\end{equation}
which is the sampling law for the reemitted packet.  Summed over all absorptions
this makes each cell radiate the correct $\kabs B_\lambda(T)$ spectrum in steady
state.  The equilibrium temperature is found by inverting the monotonic
$\kappa_{\rm P}(T)=\int\kabs(\lambda)\,B_\lambda(T)\,\dd\lambda$ against the
running absorbed power.

Both $\kappa_{\rm P}(T)$ and the correction spectrum of Eq.~\eqref{eq:bwcorr}
are known functions of wavelength, so each is integrated over its bin by an
eight-node Gauss--Legendre rule in $\ln\lambda$, with $\kabs$ read at every node
rather than once at the bin center.  The tables are built once at setup, so the
cost is irrelevant.  The reemitted wavelength is then continuous inside the
chosen bin, by the same power-law inversion as
Eq.~\eqref{eq:powerlawbin}; leaving it at the bin center would freeze the
reemitted packet's $s_{\rm ext}$ and albedo, which for the astrodust mixture
differ from their bin-averaged values by of order a percent.  The method is approximate (equilibrium only, mixture
opacity, no PAH/stochastic features) but cheap, and the emergent dust light
lands directly in the observer \code{Scattered} SED.  Over $\tau_V=0.1\ldots20$
its absorbed luminosity agrees with the Lucy method to $<2\%$.

%=============================================================================
\section{Modified Random Walk}\label{sec:mrw}
%=============================================================================

In cells that are very optically thick to \emph{scattering}, a packet would take
an enormous number of tiny steps.  The Modified Random Walk (Min et al.~2009;
Robitaille~2010) replaces a whole diffusive excursion with one step.  A sphere of
radius $R_0$ (the distance to the nearest cell wall) is inscribed at the packet
position; MRW triggers when its scattering optical thickness exceeds a threshold,
$\ksca R_0 > \gamma$ (\pp{mrw\_gamma}, default 2).  The dimensionless first-passage
variable $y$ --- the probability that the diffusion has not yet reached the sphere
surface --- has cumulative distribution
\begin{equation}
  P(y) = 2\sum_{n=1}^{\infty}(-1)^{n+1}\,y^{\,n^2},\qquad y\in[0,1],
\end{equation}
which MoCafe tabulates and inverts once at setup.  A draw of $y$ fixes the total
\emph{scattering} path length travelled inside the sphere,
\begin{equation}
  ct = -\ln y \;\cdot\; 3\,\ksca \left(\frac{R_0}{\pi}\right)^{\!2},
\end{equation}
after which the packet is moved to a random point on the sphere with a fresh
isotropic direction.  MRW diffuses with the \emph{scattering} opacity; the
absorption along the path is applied as a weight decay
$w\leftarrow w\,e^{-\kabs\,ct}$, and the corresponding absorbed energy is
deposited into the radiation-field tally as the path integral
$\int_0^{ct} w\,L_{\rm p}\,e^{-\kabs\ell}\,\dd\ell$.  Russian roulette
terminates high-albedo, high-$\tau$ excursions.  Because realistic dust is
genuinely absorbing, the diffusion is usually terminated by absorption before
MRW helps; its speed-up is large only for near-pure-scattering, very thick cells,
so MRW is off by default and adds no measurable overhead when it rarely triggers.

%=============================================================================
\section{Source (galaxy) models}\label{sec:galaxy}
%=============================================================================

A stellar component emits packets from a spatial distribution; multiple
components (\pp{nsource}) are selected in proportion to luminosity, and each may
have its own spectrum and geometry (bulge $+$ disk gives the bulge-to-total
ratio).

\subsection{Exponential disk}
A razor-thin exponential disk has surface density
$\Sigma(r)\propto e^{-r/r_s}$, so the probability of emitting from cylindrical
radius $r$ is
\begin{equation}
  p(r)\,\dd r \propto 2\pi r\,\Sigma(r)\,\dd r \propto r\,e^{-r/r_s}\,\dd r,
\end{equation}
the $\mathrm{Gamma}(2)$ (double-exponential) law.  It is sampled \emph{exactly}
as the sum of two exponential deviates, $r=-r_s(\ln\xi_1+\ln\xi_2)$
(\code{rand\_r1exp}), avoiding any inverse-CDF search.  The azimuth is uniform
(modified by spiral arms, below).

\subsection{Vertical structure}
The vertical profile is exponential, $\rho\propto e^{-|z|/z_s}$
(\code{rand\_zexp}), or the self-gravitating isothermal-sheet form
$\rho\propto\mathrm{sech}^2(z/z_s)$ (\code{rand\_sech2}), selected by
\pp{density\_zprofile}.  The same options apply to the dust disk, whose density
is the double exponential $\exp(-|z|/z_s)\exp(-r/r_s)$.

\subsection{Spiral arms}
Log-spiral arms multiply the disk by
\begin{equation}
  1 + A\,\sin\!\Big[m\Big(\frac{\ln r}{\tan p}-\phi\Big)\Big],
\end{equation}
with $m$ arms, pitch angle $p$, and amplitude $A$ (\pp{spiral\_m},
\pp{spiral\_pitch}, \pp{spiral\_amp}).  For \emph{sources} the azimuth $\phi$ is
drawn by rejection against this factor (accept if
$(1+A)\,\xi\le 1+A\sin[\dots]$); for the \emph{dust} density the cell opacity is
multiplied by the same factor.  Seen face-on, the absorbed-power (or
dust-emission) map traces the arms, validated in
\file{examples/galaxy/spiral\_faceon.in}.

\subsection{Sersic bulge: deprojection and alias sampling}\label{sec:sersic}
The projected Sersic surface-brightness profile is
$I(R)=I_0\exp[-b_m(R/R_e)^{1/m}]$ with index $m$ (\pp{src\_sersic\_index}),
effective radius $R_e$ (\pp{src\_reff}), and $b_m$ from Ciotti \& Bertin~(1999)
so that $R_e$ encloses half the projected light.  MoCafe emits from the
\emph{three-dimensional} luminosity density $\nu(r)$, recovered from $I(R)$ by the
inverse Abel transform,
\begin{equation}
  \nu(r) = -\frac{1}{\pi}\int_r^\infty \frac{\dd I}{\dd R}\,
             \frac{\dd R}{\sqrt{R^2-r^2}} .
\end{equation}
The deprojected cumulative luminosity within spherical radius $r$ is
\begin{equation}
  \mathcal L(<r) = \gamma\!\big(2m{+}1,\,b_m r^{1/m}\big)
     + \frac{2}{\pi}\frac{b_m^{\,2m+1}}{m\,\Gamma(2m{+}1)}\,r^2
       \!\int_1^\infty\!\! x\,f(x)\,e^{-b_m(rx)^{1/m}}(rx)^{1/m}\,\dd x,
\end{equation}
computed once on a logarithmic radius grid (\code{sersic\_cumulative\_3D}, with
the incomplete gamma function $\gamma$).  The radius is then drawn with the
\emph{alias} method (\S\ref{sec:alias}) rather than an inverse-CDF search: the
bin probabilities (differences of the cumulative) build an alias table, and
\code{rand\_alias\_linear} returns a continuous radius by linear interpolation
within the selected bin.  The point is placed isotropically and flattened in $z$
by the axial ratio (\pp{src\_axial\_ratio}).  The alias-sampled distribution
reproduces the analytic cumulative to $\sim10^{-4}$.

\subsection{Boxy / bar / x-bar bulges}
Generalized-ellipsoid bulges use the generalized radius
$r=(|x|^n+|y|^n+|z|^n)^{1/n}$ with boxiness $n$ (\pp{src\_boxiness}; $n=2$ is an
ellipsoid, $n>2$ boxy) and a power-law (\code{boxy}), Gaussian (\code{bar}), or
peanut (\code{xbar}) radial profile.  These are drawn with a slice-within-Gibbs
sampler (\file{random\_bulge.f90}): an auxiliary height is drawn uniformly under
the current density value, and each coordinate is then updated uniformly on the
resulting horizontal slice.  Consecutive draws are correlated (the sampler
explores by conditional updates), which is adequate for a smooth bulge but
statistically weaker than the independent alias-sampled Sersic.

%=============================================================================
\section{Random sampling}\label{sec:alias}
%=============================================================================

\subsection{Inversion and rejection}
Two elementary methods recur above.  \emph{Inversion} solves $P(x)=\xi$ for the
cumulative $P$; it is exact and consumes one deviate, used for the free path
(Eq.~\ref{eq:freepath}), the Henyey--Greenstein cosine, and the exponential
disk.  \emph{Rejection} proposes $x$ from an envelope and accepts with
probability proportional to the target density; it needs no cumulative and is
used for the spiral azimuth and inside the Gibbs bulge sampler, at the cost of
occasional rejected proposals.

\subsection{Alias method}
Walker's alias method samples an index from a discrete distribution $\{p_k\}$
($k=1\ldots K$) in $\mathcal O(1)$ time after an $\mathcal O(K)$ setup.  The
setup (\code{random\_alias\_setup}, Vose's construction) rewrites the
distribution as a set of $K$ equal-probability ``columns,'' each holding at most
two outcomes: it stores an acceptance probability $q_k$ and an alias $A_k$ for
every bin.  A draw picks a uniform bin $k$ and a uniform $\eta\in[0,1)$, and
returns $k$ if $\eta<q_k$, else $A_k$.  MoCafe uses it for wavelength sampling
(\S\ref{sec:sed}), the Mueller $S_{11}$ angle, and the Sersic radius --- anywhere
a fixed discrete distribution is sampled many times.

\subsection{Continuous piecewise-linear alias}
For a continuous variable tabulated at nodes $x_i$ with density $p_i$, the
probability of the interval $[x_i,x_{i+1}]$ is the trapezoid area
$\tfrac12(p_i+p_{i+1})(x_{i+1}-x_i)$.  An alias table over these interval
probabilities selects a bin in $\mathcal O(1)$; within the bin the value is
returned by inverting the linear density (\code{rand\_alias\_linear}), giving a
continuous, independent draw.  This is used for the Sersic radius and for
external-illumination angles.

%=============================================================================
\section{HEALPix all-sky interior observer}\label{sec:allsky}
%=============================================================================

For an observer \emph{inside} the medium (the Milky-Way case) there is no single
image plane.  Each scattering, dust-reemission, and direct event peels toward the
observer point, and the expected contribution (Eq.~\ref{eq:peel}, with the
$1/d_o^2$ replaced by the pixel solid angle) is binned into the HEALPix pixel of
the arrival direction $\nhat$.  HEALPix (G\'orski et al.~2005) tiles the sphere
into $n_{\rm pix}=12\,n_{\rm side}^2$ equal-area pixels of solid angle
$\Omega=4\pi/n_{\rm pix}$; the RING scheme is used, and the pixel index is found
by \code{vec2pix}.  The optical depth to the observer is integrated along the
matching ray (Cartesian or AMR).  The output carries the standard HEALPix
keywords so it reads directly with \code{healpy}.  Because each event maps to the
single pixel of its \emph{source} direction, a diffuse map needs many events
distributed over the sky; concentrated (inner-disk) far-infrared emission is
sparse at high $n_{\rm side}$ and benefits from a finer spatial grid or
smoothing.

%=============================================================================
\section{Energy conservation and validation}
%=============================================================================

Energy conservation is the main internal check and holds by construction:
\begin{itemize}
\item The total stellar luminosity equals the escaped starlight plus the
      absorbed (hence re-emitted) power.
\item The dust solvers normalize the emission to the absorbed power cell by
      cell, so the integrated intrinsic dust luminosity equals \code{L\_ABS} to
      machine precision, and the emergent SED integrated over $4\pi d^2$ recovers
      the same value to Monte-Carlo accuracy.
\item The $\code{\_jlam}$ output cross-checks the path-length-tallied absorbed
      power against an independent event counter.
\end{itemize}
Against external references, the \textsc{SEDust} single-cell emission tracks the
Camps et al.~(2015) seven-code stochastic-heating benchmark (which includes
\textsc{SKIRT} and \textsc{DIRTY}) to a 2.5--4.4\% median relative difference over
$U=10^{-2}\ldots10^{4}$, the Lucy and Bjorkman \& Wood absorbed luminosities
agree to $<2\%$ over $\tau_V=0.1\ldots20$, and dust emission on a uniform octree
sphere reproduces the equivalent Cartesian sphere.  These are consolidated in the
\file{examples/benchmarks/validate.py} regression harness.

%=============================================================================
\section{Summary of key numerical parameters}
%=============================================================================
\begin{table}[htb]
\centering
\caption{Parameters that control the algorithms in this note.}
\begin{tabular}{lll}
\toprule
Parameter & Meaning & Default \\
\midrule
\pp{use\_reduced\_wgt}  & forced first scattering + reduced weight & \code{.true.}\\
\pp{use\_stokes}        & polarized Mueller-matrix scattering      & \code{.false.}\\
\pp{lambda\_ref}        & reference wavelength for grey rescaling  & --- \\
\pp{dust\_model}   & grain model (\code{astrodust}/\code{dl07}/\code{mrn}/\code{zubko}) & \code{astrodust}\\
\pp{dust\_single\_teq}  & equilibrium-only emission (fast)         & \code{.false.}\\
\pp{dust\_fast\_table}  & $U$-interpolated emission table          & \code{.false.}\\
\pp{dust\_niter}        & self-absorption iterations               & 1 \\
\pp{dust\_tol}          & iteration convergence tolerance          & --- \\
\pp{mrw\_gamma}         & MRW trigger ($\ksca R_0>\gamma$)         & 2 \\
\pp{nside}              & HEALPix resolution (interior observer)   & --- \\
\bottomrule
\end{tabular}
\end{table}

%=============================================================================
\section*{References}
\addcontentsline{toc}{section}{References}
%=============================================================================
\begin{itemize}\small
\item Lucy, L.~B. 1999, A\&A, 344, 282 --- path-length mean-intensity estimator.
\item Bjorkman, J.~E. \& Wood, K. 2001, ApJ, 554, 615 --- immediate reemission
      with temperature correction.
\item Guhathakurta, P. \& Draine, B.~T. 1989, ApJ, 345, 230 --- stochastic
      heating, temperature-space recursion.
\item Draine, B.~T. \& Li, A. 2001, ApJ, 551, 807 --- thermal-discrete stochastic
      heating; grain enthalpies.
\item Draine, B.~T. \& Hensley, B.~S. 2021, ApJ, 909, 94 --- astrodust model.
\item Camps, P.\ et al. 2015, A\&A, 580, A87 --- stochastic-heating dust-emission
      benchmark.
\item Draine, B.~T. \& Li, A. 2007, ApJ, 657, 810 --- DL07 dust model.
\item Mathis, J.~S., Rumpl, W., \& Nordsieck, K.~H. 1977, ApJ, 217, 425 --- the
      $a^{-3.5}$ size distribution.
\item Draine, B.~T. \& Lee, H.~M. 1984, ApJ, 285, 89 --- graphite and silicate
      optical properties; the normalization of the MRN model used here.
\item Zubko, V., Dwek, E., \& Arendt, R.~G. 2004, ApJS, 152, 211 --- BARE-GR-S
      model.
\item Min, M.\ et al. 2009, A\&A, 497, 155; Robitaille, T.~P. 2010, A\&A, 520,
      A70 --- Modified Random Walk.
\item Jonsson, P. 2006, MNRAS, 372, 2 --- grey rescaling / polychromatic scan.
\item Ciotti, L. \& Bertin, G. 1999, A\&A, 352, 447 --- Sersic $b_m$.
\item Walker, A.~J. 1977, ACM TOMS, 3, 253; Vose, M.~D. 1991, IEEE TSE, 17, 972
      --- the alias method.
\item G\'orski, K.~M.\ et al. 2005, ApJ, 622, 759 --- HEALPix pixelization.
\item Seon, K.-I. 2010, PKAS, 25, 177 --- the $(a,g)$ single-run scan.
\item Amanatides, J. \& Woo, A. 1987, Eurographics~'87 --- fast voxel traversal.
\end{itemize}

\end{document}
